% Deux leaders modifient la même ligne en même temps ; chacun applique ensuite la
% valeur de l'autre, et les deux bases finissent croisées.
\documentclass[border=6pt]{standalone}
\usepackage{sd}
\begin{document}
\begin{tikzpicture}[x=1mm, y=1mm, note/.append style={font=\sffamily\small},
    etat/.style={draw=#1, fill=#1Bg, line width=0.7pt, rounded corners=2pt, minimum width=26mm, minimum height=7mm, font=\sffamily\footnotesize}]
  \foreach \x/\t in {0/leader A, 90/leader B} {
    \node[pillo=Dslate] at (\x,62) {\t};
    \draw[draw=Dslate, line width=0.6pt, dash pattern=on 2pt off 1.5pt] (\x,58) -- (\x,0);
  }
  \node[etat=Dslate] at (0,52) {solde 0, « init »};
  \node[etat=Dslate] at (90,52) {solde 0, « init »};
  \node[etat=Dblue] at (0,40) {solde 1, « A »};
  \node[etat=Dorange] at (90,40) {solde 2, « B »};
  \node[note, anchor=east, text=Dblue] at (-15,40) {UPDATE sur A};
  \node[note, anchor=west, text=Dorange] at (105,40) {UPDATE sur B};
  \draw[arr=Dblue] (13,36) -- node[lblc=Dblue, above, sloped, pos=0.22] {nouvelle ligne : 1, « A »} (77,18);
  \draw[arr=Dorange] (77,36) -- node[lblc=Dorange, above, sloped, pos=0.22] {nouvelle ligne : 2, « B »} (13,18);
  \node[etat=Dorange] at (0,12) {solde 2, « B »};
  \node[etat=Dblue] at (90,12) {solde 1, « A »};
  \node[note, anchor=east, align=right] at (-15,12) {update\_origin\_differs :\\noté, puis appliqué};
  \node[note, anchor=west, align=left] at (105,12) {update\_origin\_differs :\\noté, puis appliqué};
  \node[note, anchor=north, align=center] at (45,4) {Chaque leader envoie la ligne entière, pas l'opération ; chacun écrase sa valeur\\par celle de l'autre, et les deux bases restent différentes (100 lignes sur 100 mesurées).};
\end{tikzpicture}
\end{document}
